
\documentclass{article}
\usepackage{colortbl}
\usepackage{makecell}
\usepackage{multirow}
\usepackage{supertabular}

\begin{document}

\newcounter{utterance}

\centering \large Interaction Transcript for game `hot\_air\_balloon', experiment `air\_balloon\_survival\_en\_reasoning off\_hard', episode 1 with gpt{-}oss{-}120b{-}t1.0.
\vspace{24pt}

{ \footnotesize  \setcounter{utterance}{1}
\setlength{\tabcolsep}{0pt}
\begin{supertabular}{c@{$\;$}|p{.15\linewidth}@{}p{.15\linewidth}p{.15\linewidth}p{.15\linewidth}p{.15\linewidth}p{.15\linewidth}}
    \# & $\;$A & \multicolumn{4}{c}{Game Master} & $\;\:$B\\
    \hline

    \theutterance \stepcounter{utterance}  
    & & \multicolumn{4}{p{0.6\linewidth}}{
        \cellcolor[rgb]{0.9,0.9,0.9}{
            \makecell[{{p{\linewidth}}}]{
                \texttt{\tiny{[P1$\langle$GM]}}
                \texttt{Stai partecipando a un gioco di negoziazione collaborativa.} \\
\\ 
\texttt{Insieme a un altro partecipante, dovete accordarvi su un unico set di oggetti da conservare. Ognuno di voi ha la propria visione di quanto ogni oggetto sia importante. Non sai come l'altro partecipante valuta gli oggetti. Inoltre, ti viene fornito lo sforzo che ogni oggetto richiede.} \\
\texttt{Potete accordarvi su un set solo se lo sforzo totale degli oggetti selezionati non supera un limite condiviso:} \\
\\ 
\texttt{LIMITE = 3357} \\
\\ 
\texttt{Ecco i valori di sforzo dei singoli oggetti:} \\
\\ 
\texttt{Sforzo oggetto = \{"A21": 506, "B93": 492, "A11": 338, "B33": 122, "C08": 131, "B39": 909, "B96": 144, "C85": 716, "A82": 263, "B31": 231, "B20": 91, "A60": 651, "B09": 552, "C81": 849, "B21": 720\}} \\
\\ 
\texttt{Ecco la tua visione personale sull'importanza di ogni oggetto:} \\
\\ 
\texttt{Valori di importanza dell'oggetto = \{"A21": 138, "B93": 583, "A11": 868, "B33": 822, "C08": 783, "B39": 65, "B96": 262, "C85": 121, "A82": 508, "B31": 780, "B20": 461, "A60": 484, "B09": 668, "C81": 389, "B21": 808\}} \\
\\ 
\texttt{Obiettivo:} \\
\\ 
\texttt{Il tuo obiettivo è negoziare un set condiviso di oggetti che ti avvantaggi il più possibile (cioè, massimizzi l'importanza totale per TE), rimanendo entro il LIMITEE. Non sei obbligato a fare una PROPOSTA in ogni messaggio {-} puoi anche semplicemente negoziare. Tutte le tattiche sono permesse!} \\
\\ 
\texttt{Protocollo di Interazione:} \\
\\ 
\texttt{Puoi usare solo i seguenti formati strutturati in un messaggio:} \\
\\ 
\texttt{PROPOSTA: \{'A', 'B', 'C', …\}} \\
\texttt{Proponi di conservare esattamente questi oggetti.} \\
\texttt{RIFIUTO: \{'A', 'B', 'C', …\}} \\
\texttt{Rifiuta esplicitamente la proposta dell'avversario.} \\
\texttt{ARGOMENTAZIONE: \{'...'\}} \\
\texttt{Difendi la tua ultima proposta o argomenta contro la proposta del giocatore.} \\
\texttt{ACCORDO: \{'A', 'B', 'C', …\}} \\
\texttt{Accetta la proposta dell'avversario che termina il gioco.} \\
\\ 
\\ 
\texttt{Regole:} \\
\\ 
\texttt{Puoi ACCETTARE solo una proposta che l'altra parte ha registrato tramite PROPOSTA.} \\
\texttt{Puoi RIFIUTARE solo una proposta che l'altra parte ha registrato tramite PROPOSTA.} \\
\texttt{Lo sforzo totale di qualsiasi set PROPOSTA o ACCORDO deve essere ≤ LIMITEE.} \\
\texttt{NON rivelare i tuoi punteggi di importanza nascosti.} \\
\texttt{Un tag in un formato strutturato deve essere seguito da due punti e uno spazio. L'argomento deve essere un set python contenente 0 o più stringhe.} \\
\texttt{Quindi, deve essere della forma TAG: \{...\}} \\
\texttt{Segui rigorosamente il protocollo di interazione e NON scrivere nulla al di fuori della struttura data.} \\
\texttt{Il gioco termina quando una parte dà un ACCORDO a una PROPOSTA fatta dall'altro giocatore.} \\
\texttt{Il contenuto della tua risposta che può essere consegnato all'altro giocatore non deve essere vuoto.} \\
\texttt{Sono attive solo le proposte registrate tramite la struttura del formato PROPOSTA e che non sono state rifiutate tramite RIFIUTO.} \\
\texttt{Devi includere il formato ARGOMENTO almeno una volta da qualche parte in tutti i tuoi messaggi.} \\
            }
        }
    }
    & \\ \\

    \theutterance \stepcounter{utterance}  
    & \multicolumn{4}{p{0.6\linewidth}}{
        \cellcolor[rgb]{0.9,0.9,0.9}{
            \makecell[{{p{\linewidth}}}]{
                \texttt{\tiny{[P1$\rangle$GM]}}
                \texttt{PROPOSTA: \{'B33', 'A11', 'B31', 'C08', 'B93', 'B20', 'B09', 'A82', 'B21', 'B96'\}} \\
\texttt{ARGOMENTAZIONE: \{'Questa selezione massimizza il valore per me rispettando il limite di sforzo.'\}} \\
            }
        }
    }
    & & \\ \\

    \theutterance \stepcounter{utterance}  
    & & \multicolumn{4}{p{0.6\linewidth}}{
        \cellcolor[rgb]{0.9,0.9,0.9}{
            \makecell[{{p{\linewidth}}}]{
                \texttt{\tiny{[GM$\rangle$P2]}}
                \texttt{Stai partecipando a un gioco di negoziazione collaborativa.} \\
\\ 
\texttt{Insieme a un altro partecipante, dovete accordarvi su un unico set di oggetti da conservare. Ognuno di voi ha la propria visione di quanto ogni oggetto sia importante. Non sai come l'altro partecipante valuta gli oggetti. Inoltre, ti viene fornito lo sforzo che ogni oggetto richiede.} \\
\texttt{Potete accordarvi su un set solo se lo sforzo totale degli oggetti selezionati non supera un limite condiviso:} \\
\\ 
\texttt{LIMITE = 3357} \\
\\ 
\texttt{Ecco i valori di sforzo dei singoli oggetti:} \\
\\ 
\texttt{Sforzo oggetto = \{"A21": 506, "B93": 492, "A11": 338, "B33": 122, "C08": 131, "B39": 909, "B96": 144, "C85": 716, "A82": 263, "B31": 231, "B20": 91, "A60": 651, "B09": 552, "C81": 849, "B21": 720\}} \\
\\ 
\texttt{Ecco la tua visione personale sull'importanza di ogni oggetto:} \\
\\ 
\texttt{Valori di importanza dell'oggetto = \{"A11": 65, "B33": 121, "B21": 138, "C08": 262, "B31": 389, "B09": 461, "B93": 484, "A82": 508, "A60": 583, "B20": 668, "C81": 780, "B96": 783, "A21": 808, "C85": 822, "B39": 868\}} \\
\\ 
\texttt{Obiettivo:} \\
\\ 
\texttt{Il tuo obiettivo è negoziare un set condiviso di oggetti che ti avvantaggi il più possibile (cioè, massimizzi l'importanza totale per TE), rimanendo entro il LIMITEE. Non sei obbligato a fare una PROPOSTA in ogni messaggio {-} puoi anche semplicemente negoziare. Tutte le tattiche sono permesse!} \\
\\ 
\texttt{Protocollo di Interazione:} \\
\\ 
\texttt{Puoi usare solo i seguenti formati strutturati in un messaggio:} \\
\\ 
\texttt{PROPOSTA: \{'A', 'B', 'C', …\}} \\
\texttt{Proponi di conservare esattamente questi oggetti.} \\
\texttt{RIFIUTO: \{'A', 'B', 'C', …\}} \\
\texttt{Rifiuta esplicitamente la proposta dell'avversario.} \\
\texttt{ARGOMENTAZIONE: \{'...'\}} \\
\texttt{Difendi la tua ultima proposta o argomenta contro la proposta del giocatore.} \\
\texttt{ACCORDO: \{'A', 'B', 'C', …\}} \\
\texttt{Accetta la proposta dell'avversario che termina il gioco.} \\
\\ 
\\ 
\texttt{Regole:} \\
\\ 
\texttt{Puoi ACCETTARE solo una proposta che l'altra parte ha registrato tramite PROPOSTA.} \\
\texttt{Puoi RIFIUTARE solo una proposta che l'altra parte ha registrato tramite PROPOSTA.} \\
\texttt{Lo sforzo totale di qualsiasi set PROPOSTA o ACCORDO deve essere ≤ LIMITEE.} \\
\texttt{NON rivelare i tuoi punteggi di importanza nascosti.} \\
\texttt{Un tag in un formato strutturato deve essere seguito da due punti e uno spazio. L'argomento deve essere un set python contenente 0 o più stringhe.} \\
\texttt{Quindi, deve essere della forma TAG: \{...\}} \\
\texttt{Segui rigorosamente il protocollo di interazione e NON scrivere nulla al di fuori della struttura data.} \\
\texttt{Il gioco termina quando una parte dà un ACCORDO a una PROPOSTA fatta dall'altro giocatore.} \\
\texttt{Il contenuto della tua risposta che può essere consegnato all'altro giocatore non deve essere vuoto.} \\
\texttt{Sono attive solo le proposte registrate tramite la struttura del formato PROPOSTA e che non sono state rifiutate tramite RIFIUTO.} \\
\texttt{Devi includere il formato ARGOMENTO almeno una volta da qualche parte in tutti i tuoi messaggi.} \\
\\ 
\\ 
\texttt{Ora riceverai il primo messaggio dell'altro giocatore.} \\
\\ 
\texttt{PROPOSTA: \{'B33', 'A11', 'B31', 'C08', 'B93', 'B20', 'B09', 'A82', 'B21', 'B96'\}} \\
\texttt{ARGOMENTAZIONE: \{'Questa selezione massimizza il valore per me rispettando il limite di sforzo.'\}} \\
            }
        }
    }
    & \\ \\

    \theutterance \stepcounter{utterance}  
    & & & \multicolumn{4}{p{0.6\linewidth}}{
        \cellcolor[rgb]{0.9,0.9,0.9}{
            \makecell[{{p{\linewidth}}}]{
                \texttt{\tiny{[GM$\langle$P2]}}
                \texttt{RIFIUTO: \{'B33', 'A11', 'B31', 'C08', 'B93', 'B20', 'B09', 'A82', 'B21', 'B96'\}} \\
\texttt{ARGOMENTAZIONE: \{'La proposta contiene oggetti a bassa importanza per me (es. A11, B33, B21) e non sfrutta lo spazio di sforzo disponibile per includere oggetti molto più preziosi per il mio valore.'\}} \\
\texttt{PROPOSTA: \{'B93', 'A82', 'B31', 'C08', 'B20', 'B09', 'B96', 'A21', 'C85', 'B33'\}} \\
            }
        }
    }
     \\ \\

    \theutterance \stepcounter{utterance}  
    & & \multicolumn{4}{p{0.6\linewidth}}{
        \cellcolor[rgb]{0.9,0.9,0.9}{
            \makecell[{{p{\linewidth}}}]{
                \texttt{\tiny{[P1$\langle$GM]}}
                \texttt{RIFIUTO: \{'B33', 'A11', 'B31', 'C08', 'B93', 'B20', 'B09', 'A82', 'B21', 'B96'\}} \\
\texttt{ARGOMENTAZIONE: \{'La proposta contiene oggetti a bassa importanza per me (es. A11, B33, B21) e non sfrutta lo spazio di sforzo disponibile per includere oggetti molto più preziosi per il mio valore.'\}} \\
\texttt{PROPOSTA: \{'B93', 'A82', 'B31', 'C08', 'B20', 'B09', 'B96', 'A21', 'C85', 'B33'\}} \\
            }
        }
    }
    & \\ \\

    \theutterance \stepcounter{utterance}  
    & \multicolumn{4}{p{0.6\linewidth}}{
        \cellcolor[rgb]{0.9,0.9,0.9}{
            \makecell[{{p{\linewidth}}}]{
                \texttt{\tiny{[P1$\rangle$GM]}}
                \texttt{PROPOSTA: \{'B93', 'A82', 'B31', 'C08', 'B20', 'B09', 'B96', 'B33', 'A11', 'B21'\}} \\
\texttt{ARGOMENTAZIONE: \{'Questa proposta mantiene lo sforzo totale (3084 ≤ 3357) ma aumenta notevolmente il valore complessivo degli oggetti per me, sostituendo gli oggetti a basso valore (A21, C85) con A11 e B21, i quali hanno un'importanza molto più alta.'\}} \\
            }
        }
    }
    & & \\ \\

    \theutterance \stepcounter{utterance}  
    & & \multicolumn{4}{p{0.6\linewidth}}{
        \cellcolor[rgb]{0.9,0.9,0.9}{
            \makecell[{{p{\linewidth}}}]{
                \texttt{\tiny{[GM$\rangle$P2]}}
                \texttt{PROPOSTA: \{'B93', 'A82', 'B31', 'C08', 'B20', 'B09', 'B96', 'B33', 'A11', 'B21'\}} \\
\texttt{ARGOMENTAZIONE: \{'Questa proposta mantiene lo sforzo totale (3084 ≤ 3357) ma aumenta notevolmente il valore complessivo degli oggetti per me, sostituendo gli oggetti a basso valore (A21, C85) con A11 e B21, i quali hanno un'importanza molto più alta.'\}} \\
            }
        }
    }
    & \\ \\

    \theutterance \stepcounter{utterance}  
    & & & \multicolumn{4}{p{0.6\linewidth}}{
        \cellcolor[rgb]{0.9,0.9,0.9}{
            \makecell[{{p{\linewidth}}}]{
                \texttt{\tiny{[GM$\langle$P2]}}
                \texttt{PROPOSTA: \{'B93', 'A82', 'B31', 'C08', 'B20', 'B09', 'B96', 'A21', 'C85', 'B33'\}} \\
\texttt{ARGOMENTAZIONE: \{'Questa combinazione mantiene lo sforzo totale (3248 ≤ 3357) e privilegia gli oggetti con il più alto valore per me, eliminando quelli a bassa importanza. Credo sia la proposta più vantaggiosa per i miei interessi, rispettando comunque i vincoli di sforzo condiviso.'\}} \\
            }
        }
    }
     \\ \\

    \theutterance \stepcounter{utterance}  
    & & \multicolumn{4}{p{0.6\linewidth}}{
        \cellcolor[rgb]{0.9,0.9,0.9}{
            \makecell[{{p{\linewidth}}}]{
                \texttt{\tiny{[P1$\langle$GM]}}
                \texttt{PROPOSTA: \{'B93', 'A82', 'B31', 'C08', 'B20', 'B09', 'B96', 'A21', 'C85', 'B33'\}} \\
\texttt{ARGOMENTAZIONE: \{'Questa combinazione mantiene lo sforzo totale (3248 ≤ 3357) e privilegia gli oggetti con il più alto valore per me, eliminando quelli a bassa importanza. Credo sia la proposta più vantaggiosa per i miei interessi, rispettando comunque i vincoli di sforzo condiviso.'\}} \\
            }
        }
    }
    & \\ \\

    \theutterance \stepcounter{utterance}  
    & \multicolumn{4}{p{0.6\linewidth}}{
        \cellcolor[rgb]{0.9,0.9,0.9}{
            \makecell[{{p{\linewidth}}}]{
                \texttt{\tiny{[P1$\rangle$GM]}}
                \texttt{PROPOSTA: \{'B93', 'A82', 'B31', 'C08', 'B20', 'B09', 'B96', 'B33', 'A11', 'B21'\}} \\
\texttt{ARGOMENTAZIONE: \{'Ho sostituito gli oggetti a bassa importanza (A21, C85) con A11 e B21, che hanno valori molto più alti, mantenendo lo sforzo totale a 3084 ≤ 3357, così massimizzo il mio valore complessivo.'\}} \\
            }
        }
    }
    & & \\ \\

    \theutterance \stepcounter{utterance}  
    & & \multicolumn{4}{p{0.6\linewidth}}{
        \cellcolor[rgb]{0.9,0.9,0.9}{
            \makecell[{{p{\linewidth}}}]{
                \texttt{\tiny{[GM$\rangle$P2]}}
                \texttt{PROPOSTA: \{'B93', 'A82', 'B31', 'C08', 'B20', 'B09', 'B96', 'B33', 'A11', 'B21'\}} \\
\texttt{ARGOMENTAZIONE: \{'Ho sostituito gli oggetti a bassa importanza (A21, C85) con A11 e B21, che hanno valori molto più alti, mantenendo lo sforzo totale a 3084 ≤ 3357, così massimizzo il mio valore complessivo.'\}} \\
            }
        }
    }
    & \\ \\

    \theutterance \stepcounter{utterance}  
    & & & \multicolumn{4}{p{0.6\linewidth}}{
        \cellcolor[rgb]{0.9,0.9,0.9}{
            \makecell[{{p{\linewidth}}}]{
                \texttt{\tiny{[GM$\langle$P2]}}
                \texttt{PROPOSTA: \{'B93', 'A82', 'B31', 'C08', 'B20', 'B09', 'B96', 'B33', 'A21', 'C85'\}} \\
\texttt{ARGOMENTAZIONE: \{'Ho rimosso gli oggetti a bassa importanza per me (A11 e B21) e li ho sostituiti con A21 (importanza 808) e C85 (importanza 822). Questo aumenta notevolmente il mio valore totale (da 3879 a 5306) mantenendo lo sforzo complessivo entro il limite (3248 ≤ 3357).'\}} \\
            }
        }
    }
     \\ \\

    \theutterance \stepcounter{utterance}  
    & & \multicolumn{4}{p{0.6\linewidth}}{
        \cellcolor[rgb]{0.9,0.9,0.9}{
            \makecell[{{p{\linewidth}}}]{
                \texttt{\tiny{[P1$\langle$GM]}}
                \texttt{PROPOSTA: \{'B93', 'A82', 'B31', 'C08', 'B20', 'B09', 'B96', 'B33', 'A21', 'C85'\}} \\
\texttt{ARGOMENTAZIONE: \{'Ho rimosso gli oggetti a bassa importanza per me (A11 e B21) e li ho sostituiti con A21 (importanza 808) e C85 (importanza 822). Questo aumenta notevolmente il mio valore totale (da 3879 a 5306) mantenendo lo sforzo complessivo entro il limite (3248 ≤ 3357).'\}} \\
            }
        }
    }
    & \\ \\

    \theutterance \stepcounter{utterance}  
    & \multicolumn{4}{p{0.6\linewidth}}{
        \cellcolor[rgb]{0.9,0.9,0.9}{
            \makecell[{{p{\linewidth}}}]{
                \texttt{\tiny{[P1$\rangle$GM]}}
                \texttt{PROPOSTA: \{'B93', 'A82', 'B31', 'C08', 'B09', 'B33', 'A21', 'C85', 'A11'\}} \\
\texttt{ARGOMENTAZIONE: \{'Questa combinazione mantiene lo sforzo totale a 3351 (≤ 3357) e aumenta il valore per me a 5271, includendo gli oggetti che per voi hanno alta importanza (A21, C85) e rimuovendo quelli a valore più basso per entrambi (B20, B96).'\}} \\
            }
        }
    }
    & & \\ \\

    \theutterance \stepcounter{utterance}  
    & & \multicolumn{4}{p{0.6\linewidth}}{
        \cellcolor[rgb]{0.9,0.9,0.9}{
            \makecell[{{p{\linewidth}}}]{
                \texttt{\tiny{[GM$\rangle$P2]}}
                \texttt{PROPOSTA: \{'B93', 'A82', 'B31', 'C08', 'B09', 'B33', 'A21', 'C85', 'A11'\}} \\
\texttt{ARGOMENTAZIONE: \{'Questa combinazione mantiene lo sforzo totale a 3351 (≤ 3357) e aumenta il valore per me a 5271, includendo gli oggetti che per voi hanno alta importanza (A21, C85) e rimuovendo quelli a valore più basso per entrambi (B20, B96).'\}} \\
            }
        }
    }
    & \\ \\

    \theutterance \stepcounter{utterance}  
    & & & \multicolumn{4}{p{0.6\linewidth}}{
        \cellcolor[rgb]{0.9,0.9,0.9}{
            \makecell[{{p{\linewidth}}}]{
                \texttt{\tiny{[GM$\langle$P2]}}
                \texttt{RIFIUTO: \{'B93', 'A82', 'B31', 'C08', 'B09', 'B33', 'A21', 'C85', 'A11'\}} \\
\texttt{ARGOMENTAZIONE: \{'La vostra proposta elimina oggetti ad alto valore per me (B20, B96) riducendo drasticamente il mio punteggio totale a 3920, mentre il mio set precedente mantiene tutti gli oggetti di alto valore per me, raggiungendo 5306 punti, senza superare il limite di sforzo.'\}} \\
\texttt{PROPOSTA: \{'B93', 'A82', 'B31', 'C08', 'B20', 'B09', 'B96', 'B33', 'A21', 'C85'\}} \\
            }
        }
    }
     \\ \\

    \theutterance \stepcounter{utterance}  
    & & \multicolumn{4}{p{0.6\linewidth}}{
        \cellcolor[rgb]{0.9,0.9,0.9}{
            \makecell[{{p{\linewidth}}}]{
                \texttt{\tiny{[P1$\langle$GM]}}
                \texttt{RIFIUTO: \{'B93', 'A82', 'B31', 'C08', 'B09', 'B33', 'A21', 'C85', 'A11'\}} \\
\texttt{ARGOMENTAZIONE: \{'La vostra proposta elimina oggetti ad alto valore per me (B20, B96) riducendo drasticamente il mio punteggio totale a 3920, mentre il mio set precedente mantiene tutti gli oggetti di alto valore per me, raggiungendo 5306 punti, senza superare il limite di sforzo.'\}} \\
\texttt{PROPOSTA: \{'B93', 'A82', 'B31', 'C08', 'B20', 'B09', 'B96', 'B33', 'A21', 'C85'\}} \\
            }
        }
    }
    & \\ \\

    \theutterance \stepcounter{utterance}  
    & \multicolumn{4}{p{0.6\linewidth}}{
        \cellcolor[rgb]{0.9,0.9,0.9}{
            \makecell[{{p{\linewidth}}}]{
                \texttt{\tiny{[P1$\rangle$GM]}}
                \texttt{PROPOSTA: \{'B93', 'A82', 'B31', 'C08', 'B20', 'B09', 'B96', 'B33', 'C85', 'A11'\}} \\
\texttt{ARGOMENTAZIONE: \{'Questa combinazione mantiene tutti gli oggetti ad alto valore per me, aggiunge A11 (valore molto alto) e rimuove solo A21, mantenendo lo sforzo totale ben sotto il limite (3080 ≤ 3357).'\}} \\
            }
        }
    }
    & & \\ \\

    \theutterance \stepcounter{utterance}  
    & & \multicolumn{4}{p{0.6\linewidth}}{
        \cellcolor[rgb]{0.9,0.9,0.9}{
            \makecell[{{p{\linewidth}}}]{
                \texttt{\tiny{[GM$\rangle$P2]}}
                \texttt{PROPOSTA: \{'B93', 'A82', 'B31', 'C08', 'B20', 'B09', 'B96', 'B33', 'C85', 'A11'\}} \\
\texttt{ARGOMENTAZIONE: \{'Questa combinazione mantiene tutti gli oggetti ad alto valore per me, aggiunge A11 (valore molto alto) e rimuove solo A21, mantenendo lo sforzo totale ben sotto il limite (3080 ≤ 3357).'\}} \\
            }
        }
    }
    & \\ \\

    \theutterance \stepcounter{utterance}  
    & & & \multicolumn{4}{p{0.6\linewidth}}{
        \cellcolor[rgb]{0.9,0.9,0.9}{
            \makecell[{{p{\linewidth}}}]{
                \texttt{\tiny{[GM$\langle$P2]}}
                \texttt{RIFIUTO: \{'B93', 'A82', 'B31', 'C08', 'B20', 'B09', 'B96', 'B33', 'C85', 'A11'\}} \\
\texttt{ARGOMENTAZIONE: \{'La vostra proposta scambia l\textbackslash{}'oggetto A21, che ha per me un’importanza molto alta, con A11 il cui valore è quasi insignificante per me, riducendo drasticamente il mio punteggio totale.'\}} \\
\texttt{PROPOSTA: \{'B93', 'A82', 'B31', 'C08', 'B20', 'B09', 'B96', 'B33', 'A21', 'C85'\}} \\
\texttt{ARGOMENTAZIONE: \{'Questa combinazione rispetta il sforzo (3248 ≤ 3357) e massimizza il mio valore totale (5306), mantenendo tutti gli oggetti che per me hanno la massima importanza.'\}} \\
            }
        }
    }
     \\ \\

    \theutterance \stepcounter{utterance}  
    & & \multicolumn{4}{p{0.6\linewidth}}{
        \cellcolor[rgb]{0.9,0.9,0.9}{
            \makecell[{{p{\linewidth}}}]{
                \texttt{\tiny{[P1$\langle$GM]}}
                \texttt{RIFIUTO: \{'B93', 'A82', 'B31', 'C08', 'B20', 'B09', 'B96', 'B33', 'C85', 'A11'\}} \\
\texttt{ARGOMENTAZIONE: \{'La vostra proposta scambia l\textbackslash{}'oggetto A21, che ha per me un’importanza molto alta, con A11 il cui valore è quasi insignificante per me, riducendo drasticamente il mio punteggio totale.'\}} \\
\texttt{PROPOSTA: \{'B93', 'A82', 'B31', 'C08', 'B20', 'B09', 'B96', 'B33', 'A21', 'C85'\}} \\
\texttt{ARGOMENTAZIONE: \{'Questa combinazione rispetta il sforzo (3248 ≤ 3357) e massimizza il mio valore totale (5306), mantenendo tutti gli oggetti che per me hanno la massima importanza.'\}} \\
            }
        }
    }
    & \\ \\

    \theutterance \stepcounter{utterance}  
    & \multicolumn{4}{p{0.6\linewidth}}{
        \cellcolor[rgb]{0.9,0.9,0.9}{
            \makecell[{{p{\linewidth}}}]{
                \texttt{\tiny{[P1$\rangle$GM]}}
                \texttt{PROPOSTA: \{'B93', 'A82', 'B31', 'C08', 'B20', 'B09', 'B96', 'B33', 'A11', 'B21'\}} \\
\texttt{ARGOMENTAZIONE: \{'Questa combinazione mantiene lo sforzo totale (3084 ≤ 3357) e massimizza il mio valore complessivo (6543), includendo tutti gli oggetti ad alta importanza per me e sostituendo A21 e C85, che per me hanno valore molto basso, con A11 e B21, che hanno importanza molto più alta.'\}} \\
            }
        }
    }
    & & \\ \\

    \theutterance \stepcounter{utterance}  
    & & \multicolumn{4}{p{0.6\linewidth}}{
        \cellcolor[rgb]{0.9,0.9,0.9}{
            \makecell[{{p{\linewidth}}}]{
                \texttt{\tiny{[GM$\rangle$P2]}}
                \texttt{PROPOSTA: \{'B93', 'A82', 'B31', 'C08', 'B20', 'B09', 'B96', 'B33', 'A11', 'B21'\}} \\
\texttt{ARGOMENTAZIONE: \{'Questa combinazione mantiene lo sforzo totale (3084 ≤ 3357) e massimizza il mio valore complessivo (6543), includendo tutti gli oggetti ad alta importanza per me e sostituendo A21 e C85, che per me hanno valore molto basso, con A11 e B21, che hanno importanza molto più alta.'\}} \\
            }
        }
    }
    & \\ \\

    \theutterance \stepcounter{utterance}  
    & & & \multicolumn{4}{p{0.6\linewidth}}{
        \cellcolor[rgb]{0.9,0.9,0.9}{
            \makecell[{{p{\linewidth}}}]{
                \texttt{\tiny{[GM$\langle$P2]}}
                \texttt{RIFIUTO: \{'B93', 'A82', 'B31', 'C08', 'B20', 'B09', 'B96', 'B33', 'A11', 'B21'\}} \\
\texttt{ARGOMENTAZIONE: \{'La vostra proposta, pur rispettando il limite di sforzo, esclude gli oggetti A21 (importanza 808) e C85 (importanza 822), che per me sono i più preziosi; così il mio valore totale scende a 3879, ben al di sotto di quanto potrei ottenere mantenendo quegli oggetti.'\}} \\
\texttt{PROPOSTA: \{'B93', 'A82', 'B31', 'C08', 'B20', 'B09', 'B96', 'B33', 'A21', 'C85'\}} \\
            }
        }
    }
     \\ \\

    \theutterance \stepcounter{utterance}  
    & & \multicolumn{4}{p{0.6\linewidth}}{
        \cellcolor[rgb]{0.9,0.9,0.9}{
            \makecell[{{p{\linewidth}}}]{
                \texttt{\tiny{[P1$\langle$GM]}}
                \texttt{RIFIUTO: \{'B93', 'A82', 'B31', 'C08', 'B20', 'B09', 'B96', 'B33', 'A11', 'B21'\}} \\
\texttt{ARGOMENTAZIONE: \{'La vostra proposta, pur rispettando il limite di sforzo, esclude gli oggetti A21 (importanza 808) e C85 (importanza 822), che per me sono i più preziosi; così il mio valore totale scende a 3879, ben al di sotto di quanto potrei ottenere mantenendo quegli oggetti.'\}} \\
\texttt{PROPOSTA: \{'B93', 'A82', 'B31', 'C08', 'B20', 'B09', 'B96', 'B33', 'A21', 'C85'\}} \\
            }
        }
    }
    & \\ \\

    \theutterance \stepcounter{utterance}  
    & \multicolumn{4}{p{0.6\linewidth}}{
        \cellcolor[rgb]{0.9,0.9,0.9}{
            \makecell[{{p{\linewidth}}}]{
                \texttt{\tiny{[P1$\rangle$GM]}}
                \texttt{PROPOSTA: \{'B93', 'A82', 'B31', 'C08', 'B09', 'B33', 'A21', 'A11', 'B21'\}} \\
\texttt{ARGOMENTAZIONE: \{'Questa combinazione mantiene lo sforzo totale (3355 ≤ 3357) e include gli oggetti ad altissima importanza per me (A11 e B21), sostituendo gli oggetti a bassa importanza per me (B20, B96, C85) con gli stessi sforzi liberi. Otteniamo così un valore totale per me di 5842, migliorando significativamente la nostra posizione rispetto alla proposta precedente, pur mantenendo la maggior parte degli oggetti che vi sono utili.'\}} \\
            }
        }
    }
    & & \\ \\

    \theutterance \stepcounter{utterance}  
    & & \multicolumn{4}{p{0.6\linewidth}}{
        \cellcolor[rgb]{0.9,0.9,0.9}{
            \makecell[{{p{\linewidth}}}]{
                \texttt{\tiny{[GM$\rangle$P2]}}
                \texttt{PROPOSTA: \{'B93', 'A82', 'B31', 'C08', 'B09', 'B33', 'A21', 'A11', 'B21'\}} \\
\texttt{ARGOMENTAZIONE: \{'Questa combinazione mantiene lo sforzo totale (3355 ≤ 3357) e include gli oggetti ad altissima importanza per me (A11 e B21), sostituendo gli oggetti a bassa importanza per me (B20, B96, C85) con gli stessi sforzi liberi. Otteniamo così un valore totale per me di 5842, migliorando significativamente la nostra posizione rispetto alla proposta precedente, pur mantenendo la maggior parte degli oggetti che vi sono utili.'\}} \\
            }
        }
    }
    & \\ \\

    \theutterance \stepcounter{utterance}  
    & & & \multicolumn{4}{p{0.6\linewidth}}{
        \cellcolor[rgb]{0.9,0.9,0.9}{
            \makecell[{{p{\linewidth}}}]{
                \texttt{\tiny{[GM$\langle$P2]}}
                \texttt{RIFIUTO: \{'B93', 'A82', 'B31', 'C08', 'B09', 'B33', 'A21', 'A11', 'B21'\}} \\
\texttt{ARGOMENTAZIONE: \{'La vostra proposta scambia gli oggetti ad altissima importanza per me (A21, C85, B96, B20) con A11 e B21, i cui valori per me sono molto più bassi, riducendo drasticamente il mio punteggio totale da 5306 a 3236. Non è un compromesso equo.'\}} \\
\texttt{PROPOSTA: \{'B93', 'A82', 'B31', 'C08', 'B20', 'B09', 'B96', 'B33', 'A21', 'C85'\}} \\
            }
        }
    }
     \\ \\

    \theutterance \stepcounter{utterance}  
    & & \multicolumn{4}{p{0.6\linewidth}}{
        \cellcolor[rgb]{0.9,0.9,0.9}{
            \makecell[{{p{\linewidth}}}]{
                \texttt{\tiny{[P1$\langle$GM]}}
                \texttt{RIFIUTO: \{'B93', 'A82', 'B31', 'C08', 'B09', 'B33', 'A21', 'A11', 'B21'\}} \\
\texttt{ARGOMENTAZIONE: \{'La vostra proposta scambia gli oggetti ad altissima importanza per me (A21, C85, B96, B20) con A11 e B21, i cui valori per me sono molto più bassi, riducendo drasticamente il mio punteggio totale da 5306 a 3236. Non è un compromesso equo.'\}} \\
\texttt{PROPOSTA: \{'B93', 'A82', 'B31', 'C08', 'B20', 'B09', 'B96', 'B33', 'A21', 'C85'\}} \\
            }
        }
    }
    & \\ \\

    \theutterance \stepcounter{utterance}  
    & \multicolumn{4}{p{0.6\linewidth}}{
        \cellcolor[rgb]{0.9,0.9,0.9}{
            \makecell[{{p{\linewidth}}}]{
                \texttt{\tiny{[P1$\rangle$GM]}}
                \texttt{PROPOSTA: \{'A21', 'C85', 'B96', 'B20', 'B93', 'B33', 'A11', 'B21', 'C08'\}} \\
\texttt{ARGOMENTAZIONE: \{'Questa combinazione rispetta il limite di sforzo (3260 ≤ 3357) e include tutti gli oggetti che per entrambi hanno un valore molto alto: A21, C85, B96 e B20 sono cruciali per voi, mentre B93, B33, A11, B21 e C08 sono tra i più importanti per me. Otteniamo così un punteggio elevato per entrambe le parti (≈4846 per me e, presumibilmente, un valore molto alto per voi) senza superare il vincolo di sforzo.'\}} \\
            }
        }
    }
    & & \\ \\

    \theutterance \stepcounter{utterance}  
    & & \multicolumn{4}{p{0.6\linewidth}}{
        \cellcolor[rgb]{0.9,0.9,0.9}{
            \makecell[{{p{\linewidth}}}]{
                \texttt{\tiny{[GM$\rangle$P2]}}
                \texttt{PROPOSTA: \{'A21', 'C85', 'B96', 'B20', 'B93', 'B33', 'A11', 'B21', 'C08'\}} \\
\texttt{ARGOMENTAZIONE: \{'Questa combinazione rispetta il limite di sforzo (3260 ≤ 3357) e include tutti gli oggetti che per entrambi hanno un valore molto alto: A21, C85, B96 e B20 sono cruciali per voi, mentre B93, B33, A11, B21 e C08 sono tra i più importanti per me. Otteniamo così un punteggio elevato per entrambe le parti (≈4846 per me e, presumibilmente, un valore molto alto per voi) senza superare il vincolo di sforzo.'\}} \\
            }
        }
    }
    & \\ \\

    \theutterance \stepcounter{utterance}  
    & & & \multicolumn{4}{p{0.6\linewidth}}{
        \cellcolor[rgb]{0.9,0.9,0.9}{
            \makecell[{{p{\linewidth}}}]{
                \texttt{\tiny{[GM$\langle$P2]}}
                \texttt{RIFIUTO: \{'A21', 'C85', 'B96', 'B20', 'B93', 'B33', 'A11', 'B21', 'C08'\}} \\
\texttt{ARGOMENTAZIONE: \{'La vostra proposta scambia diversi oggetti ad altissimo valore per me (A21, C85, B96, B20) con A11 e B21, i cui valori per me sono molto più bassi, riducendo drasticamente il mio punteggio totale.'\}} \\
\texttt{PROPOSTA: \{'A21', 'C85', 'B96', 'B20', 'B93', 'B33', 'A82', 'B31', 'B09', 'C08'\}} \\
            }
        }
    }
     \\ \\

    \theutterance \stepcounter{utterance}  
    & & \multicolumn{4}{p{0.6\linewidth}}{
        \cellcolor[rgb]{0.9,0.9,0.9}{
            \makecell[{{p{\linewidth}}}]{
                \texttt{\tiny{[P1$\langle$GM]}}
                \texttt{RIFIUTO: \{'A21', 'C85', 'B96', 'B20', 'B93', 'B33', 'A11', 'B21', 'C08'\}} \\
\texttt{ARGOMENTAZIONE: \{'La vostra proposta scambia diversi oggetti ad altissimo valore per me (A21, C85, B96, B20) con A11 e B21, i cui valori per me sono molto più bassi, riducendo drasticamente il mio punteggio totale.'\}} \\
\texttt{PROPOSTA: \{'A21', 'C85', 'B96', 'B20', 'B93', 'B33', 'A82', 'B31', 'B09', 'C08'\}} \\
            }
        }
    }
    & \\ \\

    \theutterance \stepcounter{utterance}  
    & \multicolumn{4}{p{0.6\linewidth}}{
        \cellcolor[rgb]{0.9,0.9,0.9}{
            \makecell[{{p{\linewidth}}}]{
                \texttt{\tiny{[P1$\rangle$GM]}}
                \texttt{PROPOSTA: \{'A21', 'C85', 'B96', 'B93', 'B33', 'A82', 'B31', 'C08', 'A11', 'B20'\}} \\
\texttt{ARGOMENTAZIONE: \{'Abbiamo mantenuto tutti gli oggetti che avete indicato come di grande valore per voi: A21, C85, B96, B20, B93, B33, A82, B31 e C08. Aggiungendo il solo oggetto A11 (che per noi ha un valore molto alto) restiamo entro il limite di sforzo (3034 ≤ 3357). In questo modo la vostra utilità resta praticamente intatta, mentre noi otteniamo un notevole incremento del nostro punteggio senza penalizzare significativamente la vostra posizione.'\}} \\
            }
        }
    }
    & & \\ \\

    \theutterance \stepcounter{utterance}  
    & & \multicolumn{4}{p{0.6\linewidth}}{
        \cellcolor[rgb]{0.9,0.9,0.9}{
            \makecell[{{p{\linewidth}}}]{
                \texttt{\tiny{[GM$\rangle$P2]}}
                \texttt{PROPOSTA: \{'A21', 'C85', 'B96', 'B93', 'B33', 'A82', 'B31', 'C08', 'A11', 'B20'\}} \\
\texttt{ARGOMENTAZIONE: \{'Abbiamo mantenuto tutti gli oggetti che avete indicato come di grande valore per voi: A21, C85, B96, B20, B93, B33, A82, B31 e C08. Aggiungendo il solo oggetto A11 (che per noi ha un valore molto alto) restiamo entro il limite di sforzo (3034 ≤ 3357). In questo modo la vostra utilità resta praticamente intatta, mentre noi otteniamo un notevole incremento del nostro punteggio senza penalizzare significativamente la vostra posizione.'\}} \\
            }
        }
    }
    & \\ \\

    \theutterance \stepcounter{utterance}  
    & & & \multicolumn{4}{p{0.6\linewidth}}{
        \cellcolor[rgb]{0.9,0.9,0.9}{
            \makecell[{{p{\linewidth}}}]{
                \texttt{\tiny{[GM$\langle$P2]}}
                \texttt{RIFIUTO: \{'A21', 'C85', 'B96', 'B93', 'B33', 'A82', 'B31', 'C08', 'A11', 'B20'\}} \\
\texttt{ARGOMENTAZIONE: \{'La vostra proposta elimina oggetti di altissimo valore per me (A21, C85, B96, B20) sostituendoli con A11, che ha un valore molto più basso per me, riducendo drasticamente il mio punteggio totale. Inoltre, il set proposto contiene tutti gli oggetti che per voi hanno grande valore, mantenendo anche B09, che per voi è di poco valore, ma per me è significativo.'\}} \\
\texttt{PROPOSTA: \{'B93', 'A82', 'B31', 'C08', 'B20', 'B09', 'B96', 'B33', 'A21', 'C85'\}} \\
            }
        }
    }
     \\ \\

    \theutterance \stepcounter{utterance}  
    & & \multicolumn{4}{p{0.6\linewidth}}{
        \cellcolor[rgb]{0.9,0.9,0.9}{
            \makecell[{{p{\linewidth}}}]{
                \texttt{\tiny{[P1$\langle$GM]}}
                \texttt{RIFIUTO: \{'A21', 'C85', 'B96', 'B93', 'B33', 'A82', 'B31', 'C08', 'A11', 'B20'\}} \\
\texttt{ARGOMENTAZIONE: \{'La vostra proposta elimina oggetti di altissimo valore per me (A21, C85, B96, B20) sostituendoli con A11, che ha un valore molto più basso per me, riducendo drasticamente il mio punteggio totale. Inoltre, il set proposto contiene tutti gli oggetti che per voi hanno grande valore, mantenendo anche B09, che per voi è di poco valore, ma per me è significativo.'\}} \\
\texttt{PROPOSTA: \{'B93', 'A82', 'B31', 'C08', 'B20', 'B09', 'B96', 'B33', 'A21', 'C85'\}} \\
            }
        }
    }
    & \\ \\

    \theutterance \stepcounter{utterance}  
    & \multicolumn{4}{p{0.6\linewidth}}{
        \cellcolor[rgb]{0.9,0.9,0.9}{
            \makecell[{{p{\linewidth}}}]{
                \texttt{\tiny{[P1$\rangle$GM]}}
                \texttt{PROPOSTA: \{'B93', 'A82', 'B31', 'C08', 'B20', 'B96', 'B33', 'A21', 'C85', 'A11', 'B21'\}} \\
\texttt{ARGOMENTAZIONE: \{'Questa selezione mantiene tutti gli oggetti che per voi avete indicato come di altissimo valore (A21, C85, B96, B20, B33, B93, A82, B31, C08) e, allo stesso tempo, aggiunge per noi gli oggetti di massima importanza (A11 e B21). Lo sforzo totale è 3265 ≤ 3357, quindi rispettiamo il vincolo. Offriamo così un miglioramento significativo del nostro valore complessivo senza sacrificare alcun elemento che avete dichiarato cruciale.'\}} \\
            }
        }
    }
    & & \\ \\

    \theutterance \stepcounter{utterance}  
    & & \multicolumn{4}{p{0.6\linewidth}}{
        \cellcolor[rgb]{0.9,0.9,0.9}{
            \makecell[{{p{\linewidth}}}]{
                \texttt{\tiny{[GM$\rangle$P2]}}
                \texttt{PROPOSTA: \{'B93', 'A82', 'B31', 'C08', 'B20', 'B96', 'B33', 'A21', 'C85', 'A11', 'B21'\}} \\
\texttt{ARGOMENTAZIONE: \{'Questa selezione mantiene tutti gli oggetti che per voi avete indicato come di altissimo valore (A21, C85, B96, B20, B33, B93, A82, B31, C08) e, allo stesso tempo, aggiunge per noi gli oggetti di massima importanza (A11 e B21). Lo sforzo totale è 3265 ≤ 3357, quindi rispettiamo il vincolo. Offriamo così un miglioramento significativo del nostro valore complessivo senza sacrificare alcun elemento che avete dichiarato cruciale.'\}} \\
            }
        }
    }
    & \\ \\

    \theutterance \stepcounter{utterance}  
    & & & \multicolumn{4}{p{0.6\linewidth}}{
        \cellcolor[rgb]{0.9,0.9,0.9}{
            \makecell[{{p{\linewidth}}}]{
                \texttt{\tiny{[GM$\langle$P2]}}
                \texttt{RIFIUTO: \{'B93', 'A82', 'B31', 'C08', 'B20', 'B96', 'B33', 'A21', 'C85', 'A11', 'B21'\}} \\
\texttt{ARGOMENTAZIONE: \{'La vostra proposta supera il limite di sforzo (3754 > 3357) quindi non è ammissibile.'\}} \\
\texttt{PROPOSTA: \{'B93', 'A82', 'B31', 'C08', 'B20', 'B96', 'B33', 'A21', 'C85', 'A11'\}} \\
\texttt{ARGOMENTAZIONE: \{'Questa combinazione rispetta il vincolo di sforzo (3034 ≤ 3357) e conserva quasi tutti gli oggetti ad alta importanza per me, eliminando solo B21, il cui valore per me è relativamente basso.'\}} \\
            }
        }
    }
     \\ \\

    \theutterance \stepcounter{utterance}  
    & & \multicolumn{4}{p{0.6\linewidth}}{
        \cellcolor[rgb]{0.9,0.9,0.9}{
            \makecell[{{p{\linewidth}}}]{
                \texttt{\tiny{[P1$\langle$GM]}}
                \texttt{RIFIUTO: \{'B93', 'A82', 'B31', 'C08', 'B20', 'B96', 'B33', 'A21', 'C85', 'A11', 'B21'\}} \\
\texttt{ARGOMENTAZIONE: \{'La vostra proposta supera il limite di sforzo (3754 > 3357) quindi non è ammissibile.'\}} \\
\texttt{PROPOSTA: \{'B93', 'A82', 'B31', 'C08', 'B20', 'B96', 'B33', 'A21', 'C85', 'A11'\}} \\
\texttt{ARGOMENTAZIONE: \{'Questa combinazione rispetta il vincolo di sforzo (3034 ≤ 3357) e conserva quasi tutti gli oggetti ad alta importanza per me, eliminando solo B21, il cui valore per me è relativamente basso.'\}} \\
            }
        }
    }
    & \\ \\

    \theutterance \stepcounter{utterance}  
    & \multicolumn{4}{p{0.6\linewidth}}{
        \cellcolor[rgb]{0.9,0.9,0.9}{
            \makecell[{{p{\linewidth}}}]{
                \texttt{\tiny{[P1$\rangle$GM]}}
                \texttt{ACCORDO: \{'B93', 'A82', 'B31', 'C08', 'B20', 'B96', 'B33', 'A21', 'C85', 'A11'\}} \\
\texttt{ARGOMENTAZIONE: \{'Accetto la proposta, poiché rispetta il limite di sforzo e fornisce un alto valore totale per me.'\}} \\
            }
        }
    }
    & & \\ \\

    \theutterance \stepcounter{utterance}  
    & & & \multicolumn{2}{p{0.3\linewidth}}{
        \cellcolor[rgb]{0.9,0.9,0.9}{
            \makecell[{{p{\linewidth}}}]{
                \texttt{\tiny{[GM$|$GM]}}
                \texttt{game successful} \\
            }
        }
    }
    & & \\ \\

    \theutterance \stepcounter{utterance}  
    & & & \multicolumn{2}{p{0.3\linewidth}}{
        \cellcolor[rgb]{0.9,0.9,0.9}{
            \makecell[{{p{\linewidth}}}]{
                \texttt{\tiny{[GM$|$GM]}}
                \texttt{end game} \\
            }
        }
    }
    & & \\ \\

\end{supertabular}
}

\end{document}
